\section{Cinemática de Atitude}
A atitude de uma espaçonave pode ser representada através dos ângulos de Euler e pelas matrizes de rotação e os quaternions.
Cada uma possui vantagens e limitações quanto à interpretação geométrica, à ocorrência de singularidades e ao custo computacional.

\subsection{Ângulos de Euler}
Na convenção adotada neste trabalho, utiliza-se a sequência de rotações extrínsecas ZYX, também conhecida como \emph{yaw, pitch e roll}.
Assim, a rotação em torno de $Z$ ($\psi$) é aplicada primeiro, seguida pela rotação em torno de $Y$ ($\theta$), e por fim pela rotação em torno de $X$ ($\phi$).
As matrizes (resultante e elementares) e outros detalhes foram apresentadas na seção \eqref{sec:composicao_rotacoes_sequencia}.

\subsection{Quaternions}
\label{sec:Quaternions}
Os quaternions são empregados em aplicações aeroespaciais por causa de sua eficiência computacional e por evitarem as singularidades associadas aos ângulos de Euler.
Um quaternion unitário é definido como:
\begin{equation}
\mathbf{q} =
\begin{bmatrix}
q_0 \\ q_1 \\ q_2 \\ q_3
\end{bmatrix}
=
\begin{bmatrix}
\cos\frac{\alpha}{2} \\
u_x \sin\frac{\alpha}{2} \\
u_y \sin\frac{\alpha}{2} \\
u_z \sin\frac{\alpha}{2}
\end{bmatrix},
\end{equation}

onde $\alpha$ é o ângulo de rotação e $\vec{u} = (u_x,u_y,u_z)^T$ é o eixo unitário em torno do qual a rotação ocorre.

\subsubsection{Relação entre quaternion e matriz de rotação}

Um quaternion unitário é definido como:
\begin{equation}
q = q_0 + q_1 \, \mathbf{i} + q_2 \, \mathbf{j} + q_3 \, \mathbf{k},
\label{eq:def_quaternion}
\end{equation}
onde $q_0$ é a parte escalar e $\vec{q}_v = [q_1 \; q_2 \; q_3]^T$ é a parte vetorial.  

A rotação de um vetor $\vec{v}$ pode ser escrita como:
\begin{equation}
\vec{v}' = (q_0^2 -
\vec q_v{^T} \vec{q}_v)\,\vec{v} 
+ 2(\vec{q}_v \cdot \vec{v})\,\vec{q}_v
+ 2q_0 (\vec{q}_v \times \vec{v}).
\label{eq:rodrigues_quaternion}
\end{equation}

Reescrevendo \eqref{eq:rodrigues_quaternion} em forma matricial, obtemos:
\begin{equation}
\mathbf{R}(q) = (q_0^2 - \vec{q}_v{^T} \vec{q}_v)\,\mathbf{I}_3 
+ 2 \vec{q}_v \vec{q}_v{^T}
+ 2q_0 [\vec{q}_v]_\times,
\label{eq:rot_matrix_general}
\end{equation}

onde $[\vec{q}_v]_\times$ é a matriz antissimétrica, também chamada de operador chapéu ($\hat{\vec{q}}_v$), associada ao produto vetorial:
\begin{equation}
\hat{\vec{q}}_v = [\vec{q}_v]_\times =
\begin{bmatrix}
0 & -q_3 & q_2 \\
q_3 & 0 & -q_1 \\
-q_2 & q_1 & 0
\end{bmatrix}.
\label{eq:skew_matrix}
\end{equation}

Expandindo os termos de \eqref{eq:rot_matrix_general}, obtém-se a matriz:
\begin{equation}
\mathbf{R}(q) =
\begin{bmatrix}
1 - 2(q_2^2+q_3^2) & 2(q_1q_2 - q_0q_3) & 2(q_1q_3+q_0q_2) \\
2(q_1q_2+q_0q_3)   & 1 - 2(q_1^2+q_3^2) & 2(q_2q_3 - q_0q_1) \\
2(q_1q_3 - q_0q_2) & 2(q_2q_3+q_0q_1)   & 1 - 2(q_1^2+q_2^2)
\end{bmatrix}.
\label{eq:rot_matrix_expandida}
\end{equation}

Assim, \eqref{eq:rot_matrix_expandida} mostra a forma expandida da matriz de rotação a partir de um quaternion unitário.

\subsection{Normalização e consistência numérica}
A principal restrição dos quaternions é que eles devem permanecer normalizados:
\begin{equation}
\label{eq:normalizacao_quaternion}
\lVert \mathbf{q} \rVert^2 = q_0^2 + q_1^2 + q_2^2 + q_3^2 = 1.
\end{equation}
Na prática, erros de arredondamento e integração numérica podem violar essa condição. Para preservar a consistência numérica, adota-se a normalização periódica do vetor $\mathbf{q}$ ao longo das simulações:

\begin{equation}
    \mathbf{q} \leftarrow \frac{\mathbf{q}}{\lVert \mathbf{q} \rVert}.
\end{equation}

Nesta dissertação será adotada a convenção de notação do quaternion como
\begin{equation}
q = [\,w,\,x,\,y,\,z\,]^T,
\label{eq:quat_convencao}
\end{equation}
onde $w$ representa a parte escalar e $(x,y,z)$ correspondem à parte vetorial.

\subsection{Vantagens do uso de quaternions}
Dentre as vantagens principais destacam-se: a ausência de singularidades como o \emph{gimbal lock}, a eficiência computacional, a representação compacta (quatro parâmetros em vez de nove da matriz de rotação) e a estabilidade numérica em simulações de longa duração.
